\section{Sequential Logic}

\subsection{Fundamental Building Blocks}

\begin{definition}
Unlike combinational logic, where the output depends solely on current inputs, sequential logic produces outputs based on both current and past input values. These circuits incorporate storage elements to maintain a "state".
\end{definition}

\begin{definition}
The most basic storage element consists of two inverters connected in a feedback loop. It has two stable states ($Q=1$ or $Q=0$) and a potential metastable state where outputs oscillate.
\end{definition}

\begin{definition}
A Reset-Set (R-S) latch is typically built from cross-coupled NAND gates. 
\begin{itemize}
    \item \textbf{Set ($S=0$):} Forces the output $Q$ to $1$.
    \item \textbf{Reset ($R=0$):} Forces the output $Q$ to $0$.
    \item \textbf{Forbidden ($S=0, R=0$):} This state is invalid as it breaks the $Q = \overline{Q'}$ invariant and can lead to metastability.
\end{itemize}
\end{definition}


\begin{definition}
The Gated D Latch adds a Write Enable ($WE$) signal to an R-S latch to control when data is stored. When $WE=1$, the output $Q$ follows the data input $D$; otherwise, it holds the previous value ($Q_{prev}$).
\end{definition}


\begin{important}
Latches are "transparent" while their enable signal is high, meaning any change in $D$ immediately propagates to $Q$. This is undesirable for state registers in synchronous systems because next-state values could overwrite current-state values before the cycle ends.
\end{important}

\begin{definition}
A D Flip-Flop solves the transparency problem by using a master-slave configuration of two D latches. It is \textbf{edge-triggered}, capturing the value of $D$ only at the rising edge of the clock ($0 \rightarrow 1$ transition) and maintaining it for the entire cycle.
\end{definition}


\subsection{Timing and the Clock}

\begin{definition}
\begin{itemize}
    \item \textbf{Asynchronous:} State transitions occur immediately in response to inputs with no global coordination.
    \item \textbf{Synchronous:} State transitions are synchronized to a global \textbf{clock} signal and occur at fixed units of time.
\end{itemize}
\end{definition}

\begin{definition}
A periodic signal alternating between $0$ and $1$ that triggers state transitions across all sequential elements in a system. The clock period must be long enough to accommodate the worst-case delay of the combinational logic.
\end{definition}

\begin{important}
Standard synchronous registers use rising-clock-edge triggered flip-flops. This ensures that the state changes only once per clock cycle, making system behavior predictable and easier to design than level-triggered latch-based systems.
\end{important}

\subsection{Finite State Machines (FSM)}

\begin{definition}
A discrete-time model of a stateful system consisting of a finite number of states, external inputs, external outputs, state transitions, and output logic. 
\end{definition}

\begin{theorem}
FSMs are categorized by how they generate outputs:
\begin{itemize}
    \item \textbf{Moore FSM:} Outputs depend \textbf{only} on the current state.
    \item \textbf{Mealy FSM:} Outputs depend on \textbf{both} the current state and the current inputs.
\end{itemize}
\end{theorem}

[Image comparing Moore and Mealy FSM block diagrams]

\begin{remark}
Every FSM implementation consists of three parts:
\begin{enumerate}
    \item \textbf{State Register:} Sequential circuit storing the current state.
    \item \textbf{Next State Logic:} Combinational circuit determining the next state based on current state and inputs.
    \item \textbf{Output Logic:} Combinational circuit generating system outputs.
\end{enumerate}
\end{remark}

\begin{example}
States must be represented as binary values. Common methods include:
\begin{itemize}
    \item \textbf{Binary Encoding:} Uses the minimum number of bits ($log_2(\text{states})$); minimizes flip-flops but may lead to complex combinational logic.
    \item \textbf{One-Hot Encoding:} Uses one bit per state; simplifies next-state logic but maximizes flip-flop usage.
    \item \textbf{Output Encoding:} Uses the output signals themselves as the state bits; minimizes output logic.
\end{itemize}
\end{example}

\subsection{Memory and Registers}

\begin{definition}
A structure that stores multiple bits (e.g., 4 or 32 bits) by grouping D flip-flops together with a common clock or write-enable signal.
\end{definition}

\begin{definition}
An organization of storage locations indexed by addresses. 
\begin{itemize}
    \item \textbf{Address Space:} The total number of unique locations ($2^n$ for $n$ address bits).
    \item \textbf{Addressability:} The number of bits stored at each location.
    \item \textbf{Hardware:} Uses an **Address Decoder** to select a location and a **Multiplexer** to route the selected data to the output.
\end{itemize}
\end{definition}